\section{Coherent states, moments, and the geometry of positivity}
\label{sec:foundations}

\subsection{From an ensemble to a bosonic moment state}

Let $V=\C^D$ and $X=\Pp(V)$. For a probability measure $\mu$ on $X$ define
\[
M_t(\mu)=\int_X P_z^{\otimes t}\,d\mu(z)
\quad\text{on }\Sym^t V,
\qquad n_t=\binom{D+t-1}{t}.
\]
This is a trace-one positive operator. Schur's lemma and the trace give
$M_t(\sigma)=I_{n_t}/n_t$ for the unitary-invariant measure $\sigma$.
A weighted complex projective $t$-design is exactly a measure with this
moment. Positivity of the weights is part of the definition.

The vectors $z^{\otimes t}$ form the Veronese image. The set of possible
moment states is the convex hull of their projectors, a structured subset
of all density matrices. If $f$ is a homogeneous degree-$t$ polynomial
represented in the coherent/Bombieri inner product by a vector $v_f$, then
\[
\tr\bigl(|v_f\rangle\langle v_f|M_t(\mu)\bigr)
=\int|f(z)|^2\,d\mu(z).
\]
Thus a polynomial vanishes on the support exactly when its vector lies in
the kernel of the moment state. This simple identity is the mechanism by
which algebraic geometry becomes experimentally accessible through matched
copies of a hidden pure preparation.

More general Hermitian forms test inequalities on the support. One must
distinguish a matrix that is positive on all vectors of $\Sym^t V$ from a
matrix whose symbol is positive only on coherent vectors. The latter is a
positivity problem on an algebraic realization variety. Sum-of-squares
relaxations and ideal compression exploit that distinction
\citep{Rao2024,GouveiaParriloThomas2010}.

\fig{fig01_moment_fibres}{Moment equality preserves less information than
the ensemble itself. (a) The map $[z]\mapsto(|z_0|^2,|z_1|^2,|z_2|^2)$ sends
the nine Hesse rays to three edge midpoints, with three relative phases
above each. It sends the seven noncoordinate Singer rays to the barycenter;
the three other Singer rays are the vertices. This is a probability-simplex
projection, not a drawing of complex projective incidence. Hesse weights are
$1/9$; Singer vertex weights are $1/12$ and its other weights are $3/28$.
(b) Both ensembles and Haar measure have $M_2=I_6/6$, but the rank-two
three-copy projector $Q$ defined in the ensemble-certification theorem has
expectations $0$, $5/18$, and $1/5$, respectively. All displayed expectation
values are exact identities; decimals merely label the bars.}{fig:moments}

\subsection{Positive moment repair as a finite-dimensional inverse problem}

A useful feature of coherent moments is that small moment errors can be
repaired by adding a positive density. On $\Pp^2$ let $Q_z=P_z^{\otimes2}$,
$\tau=I_6/6$, and
\[
L(A)=(A\otimes I+I\otimes A)|_{\Sym^2\C^3},\qquad
B_z=90Q_z-30L(P_z)+5I_6.
\]
For traceless Hermitian $D$ on $\Sym^2\C^3$, the coherent fourth-moment
identity gives
\begin{equation}\label{eq:repairframe}
\int\tr(DB_z)\,d\sigma=0,\qquad
\int\tr(DB_z)Q_z\,d\sigma=D,\qquad
|\tr(DB_z)|\le51\|D\|_{\mathrm{HS}}.
\end{equation}
The number $51$ is a convenient rational upper bound: the eigenvalues of
$B_z$ are $35,-25,-25,5,5,5$, whose squared sum is $2550<51^2$.
Equation~\eqref{eq:repairframe} is an explicit right inverse of the moment
map, with a positivity radius.

\begin{proposition}[Positive repair]\label{prop:repair}
Let $\delta=\|M_2(\mu)-\tau\|_{\mathrm{HS}}>0$ and
$p=(1+102\delta)^{-1}$. There is a probability $\nu$ with
$1/2\le d\nu/d\sigma\le3/2$ such that
$\mu^*=p\mu+(1-p)\nu$ is a second design. Moreover,
\[
\Ent(\mu^*\mid\sigma)
\le p\Ent(\mu\mid\sigma)+(1-p)\log(3/2).
\]
\end{proposition}
\begin{proof}
For $D=M_2(\mu)-\tau$, take
$d\nu/d\sigma=1-\tr(DB_z)/(102\delta)$.
The preceding identities prove positivity, normalization, and exact moment
cancellation. Convexity of relative entropy proves the last inequality.
\end{proof}

This is an example of a general principle: invert the finite moment frame
first, then make the inverse perturbation small enough to remain a density.
It is stronger than an arbitrary signed correction, because positivity
allows entropy estimates and actual random sampling.

For later use the full frame operator is
\begin{equation}
 \mathsf S(B)=\int\tr(BQ_z)Q_z\,d\sigma
 =\frac{B+2L(\tr_2B)+\tr(B)I_6}{90}.
 \label{an:frame-spectrum}
\end{equation}
Here an operator on $\Sym^2\C^3$ is extended by zero on the antisymmetric
subspace before taking its partial trace. Haar's fourth tensor moment
gives the formula by permutation contraction. The scalar, adjoint, and
remaining irreducible sectors have real dimensions $1,8,27$ and frame
eigenvalues $1/6,1/15,1/90$, respectively. Thus the correction uses the
representation-theoretic decomposition of the moment space, while its
norm bound gives an explicit positive density.

\subsection{Multipoles explain the time-reversal grading}

For a spin-$j$ system, operator space decomposes under rotations as
\[
\mathcal B(V_j)\simeq\bigoplus_{\ell=0}^{2j}V_\ell.
\]
The spin-coherent symbol $a(n)=\langle n|A|n\rangle$ carries the same
multipole components, with nonzero normalization factors on each degree.
Time reversal sends $n$ to $-n$. For Hermitian $A$ invariant under time
reversal, the symbol is real and satisfies $a(n)=a(-n)$; therefore all odd
spherical-harmonic degrees vanish. On a register, simultaneous time reversal
removes odd \emph{total} multipole degree.

For qubits, the only one-site degrees are $0$ and $1$. Total degree then
coincides with body level, so odd body levels vanish. For spin one, each
nonidentity site can contribute degree $1$ or $2$, and a one-body
quadrupole is time-reversal even. The grading is therefore portable across
spins when expressed in multipoles. It should not be transported as an
odd-body rule beyond qubits. Standard irreducible-tensor and spin-coherent
conventions may be found in \citet{Radcliffe1971}.

\fig{fig11_multipoles}{Time reversal acts on multipole degree. The left two
panels show the exact restrictions of $n_z$ and $(3n_z^2-1)/2$ to a disk
parameterized by $n_z$, with a common signed color scale. Antipodal reversal
changes the dipole's sign and preserves the quadrupole. The right panel lists
possible spin-one pairs of body level and total degree for at most three
active sites, truncated at total degree four. Each active site contributes
degree one or two; blue marks even total
degree and orange marks odd total degree. In particular, the one-body,
degree-two entry survives time reversal. The identity is the entry $(0,0)$.}
{fig:multipoles}

\subsection{Two positivity notions that will recur}

For matrices $A,B$, $A\preceq B$ means $B-A$ is positive semidefinite.
For completely positive maps $\Phi,\Psi$, the notation
$\Phi\le_{\mathrm{CP}}\Psi$ means that $\Psi-\Phi$ is completely positive,
equivalently that their Choi matrices satisfy the corresponding Loewner
inequality. A completely positive comparison remains valid after tensoring
with an arbitrary identity map and applying the result to any positive
environment. This is why it is the appropriate input for network bounds.

On the geometric side, a positive polynomial witness is evaluated only on
the allowed coherent vectors. On the network side, a CP-order witness must
remain positive on all ancillary extensions. The two requirements are
different, but both replace a scalar coefficient test by positivity on the
correct realization space. This distinction is central to the later
certificates.
